\section{Introduction}
\label{sec:intro}

\subsection{The fork}

Gelfand duality identifies commutative unital $C^*$-algebras with compact
Hausdorff spaces, and noncommutative geometry proceeds by dropping
commutativity while keeping the involution.  The $C^*$-identity
$\norm{a}^2=\norm{a^*a}$ makes the norm a \emph{singular-value} radius: it is
computed from $a$ and $a^*$ together, and it is insensitive to how badly $a$
fails to be normal.

There is a second, less travelled road out of the commutative case.  On a
commutative $C^*$-algebra one equally has
\begin{equation}
  \norm{a^2}=\norm{a}^2,
  \label{eq:sq-intro}
\end{equation}
and \eqref{eq:sq-intro} makes the norm an \emph{eigenvalue} radius: it says
exactly that $\norm{a}$ equals the spectral radius $r(a)$.  On a $C^*$-algebra
the two coincide precisely on the normal elements.  Dropping commutativity while
keeping \eqref{eq:sq-intro} --- rather than keeping the involution --- is the
fork taken in these notes.

Taken literally the fork is a dead end: \eqref{eq:sq-intro} together with a
spectral reality condition forces commutativity outright
(Theorem~\ref{thm:commutativity-forced}), so the ``noncommutative'' theory has
no objects.  The proposal studied here is to relax both conditions by a
\emph{commutator budget}: a quantity $\gauge(a)$, homogeneous of degree $2$,
which vanishes exactly on the centre and therefore leaves the commutative theory
untouched, but which funds a controlled failure of \eqref{eq:sq-intro} off the
centre.  The resulting objects we call \emph{perfectly bounded algebras}, or
\emph{PB-algebras}.

\subsection{What the budget buys}

The relaxed axioms read
\[
  \norm{a}^2-\norm{a^2}\le C\,\gauge(a),
  \qquad
  \bigl\lVert\,\norm{a^2}1-a^2\,\bigr\rVert\le\norm{a^2}+C\,\gauge(a),
\]
and they are genuinely two-parameter: the first measures \emph{approximate
normality}, the second \emph{approximate spectral reality}.  With the natural
gauge $\gauge_1(a)=\norm{\ad_a}^2$ and the constant $C=\tfrac14$ the class
contains $C(X,\RR)$, the quaternions $\HH$, the matrix algebras $M_n(\RR)$ and
the upper-triangular algebras $T_n(\RR)$, and it excludes
$\RR[\varepsilon]/(\varepsilon^2)$ --- the budget is zero on a commutative
algebra, so a nilpotent thickening cannot pay for its square defect.  The
theory therefore admits noncommutative points while refusing infinitesimal ones.

\subsection{Results}

The notes that this document supersedes recorded a collection of axioms,
examples and half-settled categorical questions.  The present treatment
reorganises that material around three structural theorems and corrects two
errors.

\begin{itemize}
\item \textbf{The base is functorial} (\S\ref{sec:base}).  On the centre both
  axioms collapse to the unrelaxed ones, so $\Zc(A)\cong C(X_A,\RR)$ for a
  compact Hausdorff $X_A$: every object carries a canonical classical base,
  \emph{extracted} from the algebra rather than imposed on it.  The morphism law
  --- contractive homomorphisms that do not increase the budget --- turns out to
  force $f(\Zc(A))\subseteq\Zc(B)$ automatically
  (Lemma~\ref{lem:centre-preserved}), so $A\mapsto X_A$ is a functor
  $\PB\to\CompHaus\op$.  This is the precise sense in which the category is
  ``fibred over classical geometry''.

\item \textbf{The commutative part is coreflective}
  (Theorem~\ref{thm:coreflection}).  The centre functor $A\mapsto\Zc(A)$ is
  right adjoint to the inclusion $\CommPB\hookrightarrow\PB$, and
  $\CommPB\simeq\CompHaus\op$.  Consequently every PB-algebra is a
  $C(X_A,\RR)$-algebra and so the section algebra of an upper semicontinuous
  field over its own base (\S\ref{sec:bundle}): PB-algebras are bundles of
  pointless objects over compact Hausdorff spaces.

\item \textbf{$C=\tfrac14$ is a phase transition}
  (Theorem~\ref{thm:phase-transition}).  For $C<\tfrac14$ every object satisfies
  $r(a)\ge(1-4C)\norm{a}$, hence is semisimple and contains no nonzero nilpotent
  element; at $C=\tfrac14$ the examples $M_2(\RR)$ and $T_2(\RR)$ enter
  simultaneously, and no choice of constant separates them
  (Proposition~\ref{prop:no-constant-separates}).  The value $\tfrac14$ is thus
  not a normalisation artefact but the location of the only interesting
  threshold in the family.
\end{itemize}

Two corrections to the earlier notes are recorded where they arise.  The
identification of the monomorphisms of $\PB$ with the budget-preserving
embeddings is withdrawn: injective morphisms are already monic, and the
budget-preserving embeddings are the correct notion of \emph{subobject}, not of
monomorphism (Remark~\ref{rem:mono-correction}).  The claim that the failure of
the inclusion $\Delta\hookrightarrow M_2(\RR)$ to be legal proves incompleteness
was already retracted there; we go further and \emph{answer} the open question
it left behind.  Theorem~\ref{thm:t2-cone} exhibits an explicit legal cone
$T_2(\RR)\to M_2(\RR)$ with non-scalar diagonal image, so the scalars are not
the equalizer of $\id$ and $\Ad_{\operatorname{diag}(1,-1)}$.  Whether
\emph{some} object is remains open, and \S\ref{sec:categorical} reduces that
question to the representability of an explicit functor.

\subsection{Status discipline}

Every substantive assertion carries one of four tags.  \statusproved\ means a
complete proof is given or a precise citation is supplied.  \statusnum\ means
the statement is a finite-dimensional inequality that has been checked by the
optimisation code in \texttt{verify/}, to the tolerance reported in
\S\ref{sec:status}, but not proved.  \statusconj\ and \statusopen\ mean what they
say.  The ledger in \S\ref{sec:status} collects all of them in one table; it is
the honest summary of where the theory stands, and it is intended to shrink in
the \statusopen\ column over time.
